\scanpage{152}
\mspage{68}
% archive sheet number: 151

\[
\simeq\underset I{\underleftarrow{\lim}}\ \mathrm{Hom}_{\unsure{\mathcal E|E}}(\unsure{\zeta|\xi}_i,y),
\]

ce qui prouve que $\zeta$ est une limite inductive des $\zeta_i$ dans $\mathcal E$.

% slanted
\margin{Pour \unsure{prouver} b), on note que la pl. fid. de $\alpha$ implique que $\rho$ est \unsure{essentielt|un foncteur} \unsure{surjectif} ; comme $\rho$ commute aux $\underrightarrow{\lim}$, et que tt objet de $\unsure{\mathcal F|F}$ est \deleted{\unsure{de}} \unsure{de} la forme $\underset I{\underrightarrow{\lim}}\,x_i$, les $x_i$ dans $U$, on \unsure{aura} $\rho(x)=\underset I{\unsure{\underline{\lim}|\underrightarrow{\lim}}}\,\rho(x_i)$, les $\rho(x_i)$ dans $\rho(U)$, \unsure{et} on gagne.}

\deleted{Comme} \inserted{Ceci prouve que (moyennant c))} $M\boxtimes N$ est stable par \inserted{petites} $\underrightarrow{\lim}$, donc il reste à prouver que $M\boxtimes N$ admet une petite famille strictement génératrice. Cela \unsure{résultera|résulte} du \deleted{Lemme \unsure{n°} \unsure{su-}} \inserted{b) du dit lemme,}

\inserted{puisque $S\times T$ engendre str\textsuperscript{t} $(S\times T)^\wedge$. Cela prouve en} \deleted{\unsure{tions} de la proposition, que je vais prouver à présent.} même temps \unsure{la dernière assertion|les dernières assertions} de la proposition, sous réserve de prouver \deleted{Je vais \unsure{maintenant}} la formule (4.3.1.4), \deleted{en désignant \unsure{par} \unsure{$u=s\boxtimes t$}.}

Il faut donc prouver un iso fonctoriel en $F\in\mathrm{Ob}\,M\boxtimes N$

\[
\mathrm{Hom}(\deleted{\unsure{F},}\rho(s,t),F)\simeq F(s,t),
\]

mais par adjonction on a

\[
\begin{aligned}
\mathrm{Hom}(\deleted{\unsure{F},}\rho(s,t),F)&\simeq\mathrm{Hom}(\deleted{\alpha(F),\unsure{(s,t)}}\inserted{(s,t),\alpha(F)})=\underset{\substack{\shortparallel\\F(s,t),}}{\alpha(F)(s,t)}
\end{aligned}
\]

\deleted{ok. Pour voir que $\rho(S\times T)$ est strictement génératrice dans $M\boxtimes N$, on note que $S\times T$ est strictement génératrice dans $(S\times T)^\wedge$, et que, $\alpha$ étant pl. fid., \deleted{son adjoint à g. $\rho$ est \unsure{tjs?|ess.t} surjectif, et}}

\deleted{Reste à prouver $a\Longrightarrow c$. Mais si $M\boxtimes N$ est un ps. topos, alors tt \inserted{\unsure{contr.}} foncteur sur $M\boxtimes N$ qui}

